% !TeX root = main.tex
\section{从给定偏好到长期决策准则}\label{sec:criteria}

给定偏好的多目标序贯决策，第一步是把决策者对目标结果的取舍要求转化为明确的策略比较准则。这一转化取决于四件事：偏好以什么形式表达，作用于什么评价对象，跨阶段的结果如何累积，以及环境的随机性如何处理。只给出偏好参数，还不足以完整刻画决策要求。

本章先讨论权重、效用、参考点、优先级和约束等偏好表达方式，再分析它们如何与时间累积和随机性结合，形成具体的长期决策准则。本文所说的“决策准则”不限于实值目标函数，也包括判定策略优劣的比较规则和界定策略可接受性的条件。

\subsection{基于实值函数的多目标结果评价}\label{sec:criteria-value}

设待评价的多维结果为
\begin{equation}
\mathbf{x}=\left(x_{1},\ldots,x_{m}\right),
\end{equation}
其中 $\mathbf{x}$ 可以是即时回报、累计回报或其他明确界定的评价对象。给定偏好参数 $p$ 后，可以用标量化函数或效用函数 $u_{p}(\mathbf{x})$ 对多维结果做实值评价，从而建立结果之间的排序\cite{wang2017,hayes2022}。本节只讨论偏好在多维结果空间中的表达方式；评价对象随机时，效用如何与期望、结果分布或风险度量组合，在第~\ref{sec:criteria-stochastic}~节讨论。

线性加权是最直接的多目标评价形式：
\begin{equation}
u_{\mathbf{w}}(\mathbf{x})
=
\sum_{i=1}^{m}w_{i}x_{i},
\qquad
w_{i}\geq 0,
\qquad
\sum_{i=1}^{m}w_{i}=1.
\end{equation}
在各目标的优化方向和尺度统一后，线性加权以固定的边际替代关系描述目标之间的补偿：对任意两个目标 $x_i$ 和 $x_j$，在 $w_j>0$ 时，相对取舍由权重之比 $w_i/w_j$ 决定。Marler 和 Arora 系统分析了权重、目标取值与最终折中解之间的关系，指出权重的实际偏好含义与目标的尺度和取值范围密切相关\cite{marler2010}。因此，脱离计量单位和尺度，权重并不能解释为各目标的“重要程度”。例如，保持权重不变而改变某个成本指标的计量单位，就可能改变结果之间的相对评价。归一化可以减弱尺度差异，但所用的归一化区间、参考值和基准同样会改变给定权重对应的实际取舍。用线性权重表达偏好时，应同时写明各目标的优化方向、计量尺度和归一化方式\cite{hayes2022}。

非线性效用函数可以描述随结果水平变化的取舍关系，以及目标之间的交互作用。对于可微效用 $u(\mathbf{x})$，目标之间的局部替代关系由边际效用之比决定，因而随 $\mathbf{x}$ 的位置变化。效用的具体形式决定了允许哪些补偿、补偿如何随结果水平改变。若希望评价与各目标预先规定的改善方向一致，可以对效用施加相应的单调性条件。

公平性准则是随结果水平改变取舍的一类具体效用。对经过尺度处理、可以共同比较的正结果，广义均值为
\begin{equation}
M_q(\mathbf x)=
\begin{cases}
\left(m^{-1}\sum_{i=1}^m x_i^q\right)^{1/q},&q\neq0,\\
\left(\prod_{i=1}^m x_i\right)^{1/m},&q=0.
\end{cases}
\end{equation}
$q=1$ 对应算术平均，$q=0$ 对应几何平均，$q\to-\infty$ 时趋于最差目标 $\min_i x_i$；这里用 $q$ 表示文献中的 $p$-mean 参数，以区别于本文通用的偏好记号 $p$。这一参数族把平均收益和对最差目标的关注放进同一组给定准则中，正值域和目标尺度是其定义条件\cite{kim2025portfolio}。广义 Gini 福利则先把结果从小到大排序，再用递减的正权重评价各排序位置，使较差的结果得到更大权重\cite{siddique2020fair}。它与为各目标固定权重的线性加权不同：某个目标得到多少边际关注，取决于它当前的排序。这些效用可以作用于单次结果，也可以作用于期望结果，采用哪种随机评价需要另外规定。

当决策者更容易给出希望达到的目标水平时，可以用参考点表达偏好。Wierzbicki 将愿望水平与成就标量化函数结合，以结果相对于参考目标的实现程度组织方案比较\cite{wierzbicki1982}。设参考点为
\begin{equation}
\mathbf{z}=\left(z_{1},\ldots,z_{m}\right),
\end{equation}
表示决策者期望达到的目标水平。参考点要与偏离度量结合才构成完整的评价规则，例如加权短缺形式
\begin{equation}
d_{\mathbf{w}}(\mathbf{x},\mathbf{z})
=
\sum_{i=1}^{m}
w_{i}\max\left(z_{i}-x_{i},0\right).
\end{equation}
它只惩罚未达到参考水平的部分，衡量实际结果相对于期望水平的短缺。若决策者还关心超额完成、最大短缺或目标之间的优先关系，可以相应调整偏差函数及其聚合方式。参考点提供评价的基准，结果之间的排序由参考点与偏差度量共同决定\cite{wang2017}。

另一种聚合方式是最小化最大加权偏差。对均需最大化的结果，令 $a_i(\mathbf x)=w_i(z_i-x_i)$，Tchebycheff 标量化及其平滑形式为
\begin{equation}
d_{\mathrm T}(\mathbf x)=\max_{1\leq i\leq m}a_i(\mathbf x),
\qquad
d_{\mu}(\mathbf x)=\mu\log\!\left(\sum_{i=1}^{m}\exp\!\left(\frac{a_i(\mathbf x)}{\mu}\right)\right),
\quad \mu>0.
\end{equation}
这里的偏差保留正负号，与只惩罚短缺的加权短缺函数语义不同。最大值聚合突出偏差最大的目标，平滑形式用可微近似代替最大值运算。Lin 等给出的界
\begin{equation}
d_{\mu}(\mathbf x)-\mu\log m\leq d_{\mathrm T}(\mathbf x)\leq d_{\mu}(\mathbf x)
\end{equation}
表明平滑参数控制了标量化值的近似误差\cite{lin2024smooth}。这个界不保证任意两个结果的排序不变：评价差距较小时，平滑可能改变比较结果。因此，把静态标量化用于长期结果评价时，应写明原偏差函数和近似参数；序贯递推条件在第~\ref{sec:structure}~章另行分析。

权重、效用函数和参考点并不是彼此独立的偏好类别，而是构造评价规则的不同要素：参考点可以是效用函数或距离函数中的参数，权重可以表示目标之间或偏差之间的相对取舍。比较这些形式时，关键在于它们如何定义评价基准、允许哪些目标补偿，以及由此形成怎样的结果排序。

\subsection{优先关系与可接受性要求}\label{sec:criteria-priority}

除了用实值函数做整体评价，决策者还可以直接规定目标之间的优先关系，或用满意水平和约束界定可接受结果的范围。这类表达限定了目标之间如何补偿、哪些结果可以接受，并可与上一节的实值评价组合使用。

词典序偏好为各目标规定严格的优先顺序：比较两个结果时先看最高优先级目标，只有它无法区分两个结果时，才比较下一优先级目标。在严格词典序下，低优先级目标的改善无法补偿高优先级目标的损失。有限权重的加权能否与词典序给出一致的排序，取决于可行结果集合、目标取值范围以及高优先级目标可能出现的最小差异。如果允许高优先级目标在一定范围内退让，可以引入容差，例如限定其相对最优水平的最大允许损失，再在满足该条件的结果中优化低优先级目标\cite{skalse2022}。

严格优先关系还涉及随机结果排序的公理基础。Shakerinava 等分析了不要求连续性公理的词典序期望效用表示\cite{shakerinava2025lex}。例如，若任何有正概率严重失效的方案都劣于任意安全方案，那么再低的失效概率也无法由其他收益补偿。这种偏好违反标量期望效用表示所需的连续性，有限的惩罚权重未必能表达它。具体的安全约束是否具有这种语义，由任务定义决定。

满意水平为一个或多个目标设定希望达到或最低可接受的水平，同一个满意水平可以有两种用法。一种把未达标的程度作为软性偏差，与其他目标上的收益一起评价；另一种把满意水平作为可接受性条件，只在达标的结果或策略中继续选择。参考点方法属于前者，愿望水平不可达时仍可依据结果与参考点的偏离进行比较\cite{wierzbicki1982}；约束策略优化属于后者，用长期成本或回报阈值直接限定可接受策略集合\cite{achiam2017}。同样的数值阈值，在前一种机制中参与排序，在后一种机制中界定可接受范围。

设 $F_{i}(\pi)$ 为策略 $\pi$ 在给定长期准则下对目标 $i$ 的评价，最低可接受水平为 $b_i$，则偏好约束可以写为
\begin{equation}
F_{i}(\pi)\geq b_{i}.
\end{equation}
约束的具体含义由 $F_i(\pi)$ 的定义决定。要求期望累计结果达到阈值、要求单次执行结果达到阈值、要求结果以指定概率达到阈值，对应的是三个不同的可接受策略集合。

偏好约束和系统可行性条件的来源不同：前者反映决策者对结果可接受程度的要求，后者描述设备、资源或运行机制允许的范围，尽管二者在数学形式上可能相同。安全要求可能来自物理限制、外部规范或决策者设定的风险容忍水平，它在模型中扮演哪种角色，要结合任务定义确定。

约束界定的是可接受策略集合，不一定给出集合内部的完整排序。因此，完整的决策问题还要规定可接受集合内部的选择规则；多个要求无法同时满足时，还要给出优先关系、松弛机制或其他处理方式。优先关系、满意水平和约束分别从比较次序和可接受范围两个方面补充了实值评价。

\subsection{长期评价对象与时间聚合}\label{sec:criteria-time}

序贯决策中，偏好的评价对象由多个阶段的结果共同确定。设策略 \(\pi\) 在时刻 \(t\) 产生多维回报 \(\mathbf r_{t}\)，有限时域下的折扣累计回报为
\begin{equation}
\sum_{t = 0}^{T - 1}\gamma^{t}\mathbf r_{t},\quad\quad 0 < \gamma \leq 1,
\end{equation}
时域 \(T\) 和折扣因子 \(\gamma\) 共同规定各阶段结果如何计入\cite{roijers2013}。持续运行的任务也可以采用长期平均准则，但要说明是期望时间平均、样本路径平均还是其他形式，以及相应的极限是否存在。Brázdil 等在多目标长期平均 MDP 中分别研究了期望平均结果，以及单次运行平均结果达到阈值的概率\cite{brazdil2014}：同样称为“长期平均”，比较的随机对象可以不同。

偏好与时间累积有两种基本的组合方式。第一种先评价每个阶段的多维结果，再累积阶段评价：
\begin{equation}
\sum_{t = 0}^{T - 1}{\gamma^{t}u_{p}\left( \mathbf r_{t} \right)}.
\end{equation}
第二种先累积各目标的结果，再评价累计结果：
\begin{equation}
u_{p}\left( \sum_{t = 0}^{T - 1}{\gamma^{t}\mathbf r_{t}} \right).
\end{equation}
前者表达对每个阶段结果的取舍并把阶段评价相加，后者直接评价整个过程形成的多目标结果\cite{bauerle2014,hayes2023}。固定线性函数下两种运算可以交换，一般非线性函数下则不同。例如，对每个阶段的目标失衡分别惩罚，与只对全程累计结果的失衡惩罚，允许的阶段间补偿并不一样。

有些长期要求无法由累计回报完整表达。峰值资源占用、连续服务中断时长、首次达到目标的时间，以及全过程是否违反某项要求，都依赖结果发生的顺序或特定的路径特征。这时应先定义相应的轨迹结果，再说明偏好如何评价它。

评价对象也可以是质性结果的长期统计量。Weng 在有限状态动作、奖励只具有给定全序等级的 MDP 中，以各等级的出现次数、折扣次数或平均次数构成结果向量，并在给定参考点后评价策略的期望计数\cite{weng2011ordinal}。在其偏好公理下，参考点确定了对计数向量的线性比较，从而不必给质性等级任意赋予数值间隔。这一构造假设用等级计数就能概括历史；参考点在这里表达的是等级分布或基准表现，而不是数值收益上的短缺距离。

时间折扣与目标权重在准则中承担不同的作用：折扣因子调节不同时间阶段的相对贡献，目标权重调节不同目标之间的取舍。如果不同目标有不同的时间评价要求，例如不同的折扣因子、评价时域或累计方式，应在模型中显式保留这些差异，而不是把它们统一成时间结构相同的回报向量。

归结起来，构造长期准则时，时间维度至少要明确三点：哪些阶段的结果纳入评价，各阶段结果如何聚合，偏好在时间聚合之前还是之后作用。明确了这三点，策略实际优化的长期目标才确定下来。

\subsection{随机长期结果的评价}\label{sec:criteria-stochastic}

明确偏好表达和长期结果的构造后，还要规定随机条件下的策略评价规则。状态转移、即时回报和策略执行都可能是随机的，同一策略在不同运行中会产生不同的长期结果，因此要说明效用如何作用于随机结果。以累计回报向量为评价对象时，Hayes 等将两种基本形式区分为标量化期望回报（scalarised expected returns，SER）与期望标量化回报（expected scalarised returns，ESR）\cite{hayes2022}。

设策略 $\pi$ 产生的随机累计回报为 $\mathbf{G}^{\pi}$，偏好参数为 $p$。SER 先对各目标的累计回报取期望，再对期望结果做效用评价：
\begin{equation}
J_{\mathrm{SER}}(\pi;p)
=
u_p\!\left(
\mathbb{E}_{\pi}\!\left[\mathbf{G}^{\pi}\right]
\right).
\end{equation}
ESR 先对每次执行得到的累计结果做效用评价，再对效用取期望：
\begin{equation}
J_{\mathrm{ESR}}(\pi;p)
=
\mathbb{E}_{\pi}\!\left[
u_p\!\left(\mathbf{G}^{\pi}\right)
\right].
\end{equation}

二者的差别在于效用评价与随机结果聚合的先后顺序。$u_p$ 为线性函数时两种运算可以交换；对一般非线性效用，两种准则通常给出不同的策略排序。以两个均需最大化的目标为例，设
\begin{equation}
u(x_1,x_2)=\min\{x_1,x_2\}.
\end{equation}

两种准则利用结果分布信息的方式也不同。SER 只依据期望累计回报向量评价策略，期望向量相同的策略在 SER 下评价相同；ESR 让非线性效用作用于每个可能的结果，因此对结果分布的差异有响应。这种响应的含义仍由效用函数决定：同一个非线性形式，既可以描述目标之间的取舍，也可以使评价对随机结果的波动或尾部表现敏感。使用非线性效用时，应区分它在目标权衡和随机结果评价中各起什么作用。

SER 与 ESR 描述的是效用与期望运算的组合关系。对于第~\ref{sec:criteria-time}~节中的阶段效用累积或其他轨迹评价对象，随机评价要按相应的长期结果定义进行：阶段效用累积先在单条轨迹内完成时间聚合，再对不同运行结果做统计；峰值、首次达到时间或路径事件等特征，则先从完整轨迹中提取评价量，再规定其在随机条件下的比较方式。

期望之外，随机决策要求还可以用尾部风险、达标概率或约束违反概率表达。条件风险价值（Conditional Value-at-Risk，CVaR）刻画损失分布较差尾部的平均损失，是风险敏感优化中常用的表示\cite{rockafellar2000}。多目标任务使用这类风险指标时，要写明风险作用于哪个目标、哪种目标组合或哪个长期评价量。若要求多个目标的累计结果以给定概率同时达到阈值，可以写为
\begin{equation}\label{eq:joint-chance}
\Pr_{\pi}\!\left(
G_i^{\pi}\geq b_i,\ \forall i
\right)
\geq 1-\alpha.
\end{equation}

这两类要求与 SER、ESR 的关系可以从运算顺序看出。式~\eqref{eq:joint-chance} 左侧等于 $\mathbb E_\pi[\mathbf 1\{G_i^\pi\geq b_i,\ \forall i\}]$，即以联合达标指示函数为效用的 ESR，因此它依赖每条轨迹上各目标的联合结果，而不能由期望回报向量确定。CVaR 是作用于整个损失分布的风险度量，不是对某个固定效用取期望，但同样只能由长期结果的分布计算。所以，达标概率和尾部风险与 ESR 一样，要求求解和评价保留轨迹层面的结果信息，这一点在第~\ref{sec:structure}~章讨论状态与价值表示时还会用到。

\subsection{长期决策准则及其组成}\label{sec:criteria-summary}

给定偏好下的长期决策准则由偏好表达、长期评价对象、时间聚合方式和随机性处理共同确定。对于能用实值函数评价并包含约束的情形，可以统一写为
\begin{equation}
\begin{aligned}
\max_{\pi\in\Pi_{\mathrm{phys}}}\quad
& J_p(\pi) \\
\mathrm{s.t.}\quad
& C_{p,j}(\pi)\leq 0,
\qquad j=1,\ldots,k ,
\end{aligned}
\end{equation}
其中 $\Pi_{\mathrm{phys}}$ 是满足系统可行性条件的策略集合，$J_p(\pi)$ 是偏好参数 $p$ 下的长期策略评价，$C_{p,j}(\pi)$ 是相应的可接受性约束。词典序等无法自然写成单一实值目标的偏好，则采用相应的排序关系或分层优化规则。

表~\ref{tab:criterion-definition}~将准则的组成与可接受性、可行性检查对应起来。
\begin{table}[H]
\centering
\caption{长期决策准则的组成与语义检查}
\label{tab:criterion-definition}
\begingroup
\zihao{5}\songti
\setlength{\tabcolsep}{6pt}
\renewcommand{\arraystretch}{1.18}
\begin{tabularx}{\textwidth}{@{}P{0.17\textwidth}Y Y@{}}
\toprule
组成 & 核心问题 & 典型表达 \\
\midrule
偏好表达 & 目标之间如何取舍？ & 权重、效用、参考点、词典序 \\
\addlinespace[4pt]
长期评价对象 & 评价什么结果？ & 阶段结果、累计向量、峰值、路径事件 \\
\addlinespace[4pt]
时间聚合 & 各阶段如何合并？ & 阶段效用累积、累计结果效用、长期平均 \\
\addlinespace[4pt]
随机性处理 & 随机结果如何比较？ & SER、ESR、达标概率、指定损失的尾部风险 \\
\midrule
跨项检查 & 什么策略可接受或可行？ & 期望阈值、联合达标要求、物理可行域 \\
\bottomrule
\end{tabularx}
\par\vspace{4pt}\raggedright\noindent
注：前四项构成准则，末行为可接受性与可行性检查，各项可以组合。还须明确目标方向与尺度、时域与折扣、要求来源及不可满足时的处理。SER/ESR、优先关系与期望约束参见文献\cite{hayes2022,skalse2022,achiam2017}；相同阈值作用于不同对象时，可能定义不同问题。
\par\endgroup
\end{table}


不同的数学表达是否对应同一个决策问题，要看它们是否保持原有的策略排序和可接受集合。用惩罚项代替硬约束、用有限权重近似严格优先关系、改变效用与期望的组合顺序，都可能改变原来的决策准则。

长期准则确定之后，还要分析它能否逐阶段分解，以及求解时需要保留哪些状态和价值信息。下一章围绕可加性、状态充分性、递推关系和策略形式展开讨论。
